\begin{Abstract}
Large Language Models (LLMs) inherently absorb harmful knowledge, misinformation, and personal data during pretraining on large-scale datasets sourced from the internet, and they are simultaneously exposed to backdoor attacks in which vicious triggers added during training or model editing elicit harmful outputs on specific input patterns while maintaining clean performance on normal inputs. Machine unlearning offers a principled path to remove such unwanted content without retraining from scratch. However, the same threat surface raises two distinct and largely unaddressed problems. First, legitimate watermarks used as ownership signatures share similar mechanisms to backdoors, creating a critical challenge of detecting and eliminating unknown backdoors without compromising watermark integrity. Second, existing unlearning methods are provider-centric, requiring training pipelines, curated retain datasets, and direct intervention by model service providers (MSPs), which leaves end users without any control over their own data. This thesis develops two complementary frameworks that address these problems under minimal assumptions.

The first framework, \textbf{BackFlush}, is a unified framework for backdoor detection and elimination while preserving watermarks. Existing defenses require prior knowledge of triggers or their payloads, depend on clean reference models, or sacrifice model utility without preserving the watermark. To address these limitations, BackFlush establishes two novel observations. The \textit{Backdoor Flushing Phenomenon} shows that injecting and unlearning auxiliary data eliminates pre-established backdoors. \textit{Backdoor Susceptibility Amplification} enables constant-time detection independent of vocabulary size. BackFlush then employs Rotation-based Parameter Editing (RoPE) unlearning, a technique that preserves watermarks while eliminating backdoors by rotating the embeddings toward their opposite direction through a bounded rotation loss, in contrast to gradient ascent whose unbounded gradients indiscriminately corrupt both backdoor and watermark subspaces.

The second framework addresses the user-control gap by formalizing \textit{Interactive Machine Unlearning} (IMU), a novel problem setting in which users can instruct LLMs to unlearn targeted knowledge through natural language prompts. This process occurs during inference, eliminating the need for any middleman. To solve IMU, we propose the \textbf{RePAIR} framework, which consists of a watchdog model that detects unlearning intent, a surgeon model that generates repair code, and a patient model whose weights are updated autonomously. At its core, we introduce \textbf{S}teering \textbf{T}hrough \textbf{A}ctivation \textbf{M}anipulation with \textbf{P}seudo\textbf{I}nverse \textbf{(STAMP)}, a training-free, single-sample unlearning method that redirects MLP activations toward a refusal subspace via closed-form pseudoinverse updates. A low-rank variant further reduces computational cost from $\mathcal{O}(d^3)$ to $\mathcal{O}(r^3 + r^2 \cdot d)$, achieving up to an approximately 3$\times$ speedup over training-based baselines.

Experimental validation confirms the effectiveness of both frameworks. Comprehensive evaluation of BackFlush across diverse trigger types over different architectures demonstrates an approximately 1\% Attack Success Rate (ASR) and approximately 99\% Clean Accuracy (CACC) while preserving watermarking capabilities, occupying a regime where no existing method simultaneously provides these properties alongside maintaining model utility comparable to clean baselines. For RePAIR, experiments across three tasks, namely harmful knowledge suppression, misinformation correction, and personal data erasure, demonstrate near-zero forget scores ($\mathrm{Acc}_f = 0.00$, $F\text{-}RL = 0.00$) while preserving retention ($\mathrm{Acc}_r$ up to 84.47, $R\text{-}RL$ up to 0.88), outperforming six state-of-the-art baselines and achieving near-oracle forgetting while preserving model utility.

Together, these frameworks take a step toward responsible and democratized model governance. BackFlush enables third-party models to be audited and sanitized before deployment without prior knowledge of triggers or payloads and without compromising intellectual property, while RePAIR enables transparent, on-device, user-centric control over learned knowledge, allowing end users to exercise their right to be forgotten directly through natural language at inference time.
\\
\textbf{Keywords: -} Machine unlearning, Large Language Models, backdoor detection and elimination, watermark preservation, knowledge-free defense, constant-time detection, Rotation-based Parameter Editing, interactive machine unlearning, training-free unlearning, single-sample forgetting, closed-form pseudoinverse update, low-rank approximation, harmful knowledge suppression, misinformation correction, personal data erasure, on-device user-centric model governance.
\end{Abstract}